\documentclass[conference]{IEEEtran}
\IEEEoverridecommandlockouts
\usepackage{cite}
\usepackage{amsmath,amssymb,amsfonts}
\usepackage{algorithmic}
\usepackage{graphicx}
\usepackage{textcomp}
\usepackage{xcolor}
\usepackage{booktabs}
\usepackage{multirow}
\def\BibTeX{{\rm B\kern-.05em{\sc i\kern-.025em b}\kern-.08em
    T\kern-.1667em\lower.7ex\hbox{E}\kern-.125emX}}
\begin{document}

\title{A Unified Framework for Low-Resource African Natural Language Processing:
Cross-Lingual Transfer, Embedding Augmentation, and Knowledge Distillation%
\thanks{Supported by the Data Science for Social Impact (DSFSI) Research
Laboratory, University of Pretoria.}}

\author{
\IEEEauthorblockN{1\textsuperscript{st} Omprakash Pugazhendhi}
\IEEEauthorblockA{\textit{Dept.\ of Computer Science and Engineering} \\
\textit{Vellore Institute of Technology}\\
Chennai, India \\
omprakash.2021@vitstudent.ac.in}
\and
\IEEEauthorblockN{2\textsuperscript{nd} Vukosi Marivate}
\IEEEauthorblockA{\textit{Dept.\ of Computer Science} \\
\textit{University of Pretoria}\\
Pretoria, South Africa \\
vukosi.marivate@up.ac.za}
}

\maketitle

% ---------------------------------------------------------------
\begin{abstract}
Natural language processing for African languages is held back by three
problems that compound each other: there is very little labelled training
data, the languages are highly diverse even within the same family, and
the best-performing models are far too large to run in the low-bandwidth,
limited-compute environments where these languages are most spoken.
This paper addresses all three barriers in a unified study spanning three
experiments on publicly available African-language benchmarks.

In the first experiment, we benchmark how well six model architectures
transfer knowledge from Kinyarwanda to the closely related Kirundi language
for a 12-class news classification task. We introduce a
\emph{catastrophic forgetting percentage} (CF\%) to measure how much each
model degrades on its source language after retraining on the target.
AfriBERTa achieves the best Kirundi accuracy at 88.30\% with only 5.14\%
source-language degradation.

In the second experiment, we propose MAGE (Multi-Head Attention Guided
Embeddings), a fully language-independent pipeline for tweet sentiment
analysis. MAGE augments sentence embeddings with small random perturbations,
refines them through an autoencoder, and classifies with a four-head
attention ensemble. On combined AfriSenti data for Kinyarwanda, Swahili,
and Xitsonga, MAGE with a denoising autoencoder improves accuracy by
3.64 percentage points over a logistic regression baseline.

In the third experiment, we compress AfroXLMR-Large (559.9M parameters,
2.1~GB) into AfroXLMR-Comet (68.9M parameters, 263~MB) by simultaneously
training the student to match the teacher's output probabilities and
internal attention patterns. AfroXLMR-Comet achieves 14.0~ms inference
latency --- 21$\times$ faster than the teacher --- and retains 89.2\%
of teacher F1 across five African languages.
\end{abstract}

\begin{IEEEkeywords}
low-resource NLP, African languages, cross-lingual transfer, knowledge
distillation, data augmentation, Kinyarwanda, Kirundi, sentiment analysis
\end{IEEEkeywords}

% ---------------------------------------------------------------
\section{Introduction}

Africa is home to more than 2,000 living languages, yet fewer than twenty
appear in mainstream NLP benchmarks \cite{b_adelani_afrisenti}.
This gap means that hundreds of millions of speakers of languages such as
Kinyarwanda, Kirundi, Swahili, Hausa, Igbo, and Yoruba are largely
excluded from AI-powered services that depend on language understanding ---
search, content moderation, health information delivery, and educational
tools among them.

Three structural barriers drive this exclusion.
The first is \emph{data scarcity}: even the largest available labelled
corpora for Bantu languages contain only tens of thousands of sentences,
compared with billions for English.
Training a model from scratch on such data is impractical.
The second is \emph{linguistic diversity}: even closely related languages
like Kinyarwanda and Kirundi share only partial vocabulary, so transferring
a model across them is not straightforward.
The third is \emph{computational asymmetry}: state-of-the-art multilingual
models contain hundreds of millions of parameters, making them impractical
for deployment in Sub-Saharan Africa where GPU infrastructure is scarce
and connectivity is intermittent.

This paper presents three experiments that directly address each barrier.

\textbf{Contribution 1.}
We systematically benchmark zero-shot cross-lingual transfer from
Kinyarwanda to Kirundi using the KINNEWS/KIRNEWS news corpus \cite{b_kinnews}.
Six architectures are compared: mBERT \cite{b_devlin}, AfriBERTa
\cite{b_ogueji}, BantuBERTa \cite{b_dossou}, BiGRU \cite{b_cho},
TextCNN \cite{b_kim}, and CharCNN \cite{b_zhang}. We define a
\emph{catastrophic forgetting percentage} (CF\%) as the fraction of
source-language accuracy lost after target-language fine-tuning.

\textbf{Contribution 2.}
We implement MAGE, a three-stage pipeline for low-resource sentiment
classification using AfriSenti data \cite{b_adelani_afrisenti}.
The pipeline applies Language-Independent Data Augmentation (LiDA) in
the embedding space, refines augmented embeddings with an autoencoder
(DAE or VAE), and classifies using a multi-head attention ensemble.
The entire pipeline is language-agnostic: it requires no dictionaries,
translation systems, or language-specific lexical resources.

\textbf{Contribution 3.}
We implement AfroXLMR-Comet, a compact student model distilled from
AfroXLMR-Large \cite{b_alabi} using a hybrid loss that combines KL
divergence on soft output distributions and MSE on internal attention maps.
The student is then fine-tuned on five African languages and evaluated
on AfriSenti.

The paper is organised as follows. Section~\ref{sec:related} reviews
related work. Section~\ref{sec:data} describes the datasets.
Section~\ref{sec:method} details all three methodologies.
Section~\ref{sec:results} reports experimental results.
Section~\ref{sec:discussion} interprets the findings.
Section~\ref{sec:conclusion} concludes.

% ---------------------------------------------------------------
\section{Related Work}\label{sec:related}

\subsection{Multilingual Pretrained Language Models}
mBERT \cite{b_devlin} and XLM-R \cite{b_conneau} demonstrated that a
single Transformer pretrained on 100+ languages can transfer zero-shot
to unseen target languages through shared sub-word vocabulary.
AfroXLMR \cite{b_alabi} further adapted XLM-R via continued pretraining
on 17 African languages. AfriBERTa \cite{b_ogueji} and BantuBERTa
\cite{b_dossou} are purpose-built models trained exclusively on African
corpora, enabling stronger within-family transfer for Bantu languages.

\subsection{Cross-Lingual Transfer and Catastrophic Forgetting}
Zero-shot cross-lingual transfer exploits shared sub-word tokens across
related languages \cite{b_conneau}. A known failure mode is catastrophic
forgetting \cite{b_mccloskey}: when fine-tuned on a new language, a model
can lose accuracy on the original. Elastic Weight Consolidation
\cite{b_kirkpatrick} combats this by penalising changes to the most
task-critical parameters. Our CF\% metric quantifies this degradation
cleanly and consistently across all six architectures, without modifying
the training procedure.

\subsection{Data Augmentation for Low-Resource NLP}
Easy Data Augmentation (EDA) \cite{b_wei} applies synonym replacement and
random word operations, but requires language-specific lexical resources
unavailable for most African languages. Back-translation
\cite{b_sennrich} generates extra training examples via machine
translation, but reliable translation systems do not exist for most
Bantu languages. LiDA \cite{b_mage} bypasses both limitations by adding
small random perturbations directly to sentence embeddings rather than
manipulating text, making it fully language-agnostic. Denoising
autoencoders \cite{b_vincent} and variational autoencoders \cite{b_kingma}
then provide principled ways to regularise the perturbed representations.

\subsection{Knowledge Distillation}
Hinton et al.\ \cite{b_hinton} introduced knowledge distillation: training
a student to match a teacher's soft output distributions rather than hard
ground-truth labels, transferring the teacher's confidence patterns
across classes. DistilBERT \cite{b_sanh} applied this to BERT, retaining
97\% of teacher performance at 40\% fewer parameters. TinyBERT
\cite{b_jiao} added intermediate layer matching including attention-map
alignment. Our distillation approach follows the attention transfer idea
from \cite{b_zagoruyko}, using mean-pooled attention matrices with a
learnable projection to handle the dimensional gap between a 16-head
teacher and an 8-head student.

% ---------------------------------------------------------------
\section{Datasets}\label{sec:data}

\subsection{KINNEWS and KIRNEWS}
KINNEWS and KIRNEWS \cite{b_kinnews} contain online news articles in
Kinyarwanda and Kirundi across 14 topical categories (Politics, Sport,
Economy, Health, Entertainment, History, Technology, Tourism, Culture,
Fashion, Religion, Environment, Education, Relationships). Following
\cite{b_xlt2024}, we remove the two smallest classes (Fashion and
Tourism) and remap the remaining 12 labels to a contiguous zero-indexed
range. Article titles and body text are concatenated with a
\texttt{[SEP]} separator before tokenisation. Dataset sizes after
preprocessing are shown in Table~\ref{tab:data_news}.

\begin{table}[htbp]
\caption{News Classification Datasets (after label filtering)}
\label{tab:data_news}
\begin{center}
\begin{tabular}{lrrc}
\toprule
\textbf{Language} & \textbf{Train} & \textbf{Test} & \textbf{Classes} \\
\midrule
Kinyarwanda & 17,014 & 4,254 & 12 \\
Kirundi     &  3,690 &   922 & 12 \\
\bottomrule
\end{tabular}
\end{center}
\end{table}

\subsection{AfriSenti}
AfriSenti \cite{b_adelani_afrisenti} provides tweet-level sentiment
annotations (positive, negative, neutral) for eleven African languages
from SemEval-2023 Shared Task 12. We use three languages for the MAGE
augmentation pipeline (Section~\ref{sec:mage}) and five for the
distillation fine-tuning stage (Section~\ref{sec:distill}). Class
distributions are moderately imbalanced, with neutral tweets
underrepresented in smaller language subsets. Dataset splits are shown
in Table~\ref{tab:data_sent}.

\begin{table}[htbp]
\caption{AfriSenti Sentiment Dataset Splits}
\label{tab:data_sent}
\begin{center}
\small
\begin{tabular}{llrr}
\toprule
\textbf{Language} & \textbf{Experiment} & \textbf{Train} & \textbf{Test} \\
\midrule
Kinyarwanda & MAGE + Distil & 3,302  & 1,026 \\
Swahili     & MAGE + Distil & 1,810  &   748 \\
Xitsonga    & MAGE only     & 1,006  &   252 \\
Hausa       & Distil only   & 14,172 & 5,303 \\
Igbo        & Distil only   & 10,192 & 3,682 \\
Yoruba      & Distil only   &  8,522 & 4,515 \\
\bottomrule
\end{tabular}
\end{center}
\end{table}

% ---------------------------------------------------------------
\section{Methodology}\label{sec:method}

\subsection{Cross-Lingual Transfer Benchmark}\label{sec:xlt}

\subsubsection{Transformer Models}
We fine-tune three pretrained transformer models on the Kinyarwanda
training set and evaluate on Kirundi under two conditions: (1) zero-shot,
with no Kirundi exposure at all, and (2) after further fine-tuning on the
Kirundi training set. All models are loaded as 12-class sequence
classifiers via HuggingFace Transformers \cite{b_wolf}.

The shared training configuration is: batch size 32, learning rate
$5 \times 10^{-5}$, weight decay 0.01, 500 warm-up steps, and maximum
input length 128 tokens. mBERT and BantuBERTa train for 8 epochs;
AfriBERTa trains for 25 epochs to allow its smaller, Africa-specific
vocabulary to fully converge. Checkpoints are selected by best macro F1.

\subsubsection{Traditional Neural Models}
Three non-transformer architectures are evaluated alongside the
transformers. They are far smaller and cheaper to train, making them
useful baselines for settings with very limited compute.

\textbf{BiGRU.} A two-layer bidirectional GRU \cite{b_cho} over
50-dimensional skip-gram Word2Vec embeddings \cite{b_mikolov} trained
from scratch on the combined Kinyarwanda and Kirundi corpora (window 5,
min count 5, hierarchical softmax). Hidden dimension 256; dropout 0.5;
10 training epochs with Adam.

\textbf{TextCNN.} The architecture from \cite{b_kim}: parallel
convolutions with filter widths $\{3, 4, 5\}$ (150 feature maps each)
over the same Word2Vec embeddings, max-over-time pooling, dropout 0.5,
linear classifier. Trained for 8 epochs.

\textbf{CharCNN.} The architecture from \cite{b_zhang}: 68-character
vocabulary, sequences padded or truncated to 1,500 characters, six
convolutional layers (256 filters, kernel 7), three max-pooling layers,
two fully connected layers (1,024 units). Trained with SGD at
learning rate 0.01 for 30 epochs.

\subsubsection{Catastrophic Forgetting Metric}
After a model has been trained on Kinyarwanda and evaluated on Kirundi,
we fine-tune it further on Kirundi and then re-evaluate it on Kinyarwanda.
The catastrophic forgetting percentage measures how much Kinyarwanda
accuracy was lost during that Kirundi retraining:
\begin{equation}
    \text{CF\%} = \frac{
      \text{Acc}_{\text{KIN}}^{\text{before}} -
      \text{Acc}_{\text{KIN}}^{\text{after}}
    }{
      \text{Acc}_{\text{KIN}}^{\text{before}}
    } \times 100
    \label{eq:cf}
\end{equation}
``Before'' is Kinyarwanda accuracy after training only on Kinyarwanda;
``after'' is Kinyarwanda accuracy once the same model has been further
fine-tuned on Kirundi. A CF\% of zero means the model kept all its
source-language knowledge. As a concrete example: a model that scores
85\% before and 80\% after Kirundi fine-tuning has
CF\% $= (85 - 80)/85 \times 100 = 5.9\%$, meaning it lost about 6\%
of its original capability.

\subsection{MAGE: Augmentation and Attention Ensemble}\label{sec:mage}

MAGE is a four-stage pipeline: choose a sentence encoder, augment
embeddings with LiDA, refine with an autoencoder, and classify with a
multi-head attention ensemble. Each stage can be ablated independently.

\subsubsection{Sentence Encoding and Encoder Selection}
Four candidate encoders are compared by extracting mean-pooled
last-hidden-state embeddings and classifying with logistic regression
on the combined AfriSenti test set. Results are shown in
Table~\ref{tab:embed}. AfriBERTa achieves the highest F1 (0.580)
and is used for all subsequent MAGE stages.

\begin{table}[htbp]
\caption{Sentence Encoder Comparison on Combined AfriSenti Test Set}
\label{tab:embed}
\begin{center}
\begin{tabular}{lcccc}
\toprule
\textbf{Encoder} & \textbf{Acc} & \textbf{F1} & \textbf{Prec} & \textbf{Recall} \\
\midrule
mBERT          & 0.483 & 0.425 & 0.464 & 0.455 \\
Afro XLM-R     & 0.520 & 0.522 & 0.528 & 0.549 \\
BantuBERTa     & 0.526 & 0.513 & 0.520 & 0.511 \\
AfriBERTa      & \textbf{0.569} & \textbf{0.580} & \textbf{0.578} & \textbf{0.560} \\
\bottomrule
\end{tabular}
\end{center}
\end{table}

\subsubsection{Language-Independent Data Augmentation (LiDA)}
Given an AfriBERTa sentence embedding $\mathbf{e} \in \mathbb{R}^{768}$,
LiDA generates an augmented copy by adding uniform random noise
independently to every dimension:
\begin{equation}
    \tilde{\mathbf{e}} = \mathbf{e} + \boldsymbol{\epsilon},
    \quad
    \boldsymbol{\epsilon} \sim \mathcal{U}(r_{\min},\, r_{\max})^{768}
    \label{eq:lida}
\end{equation}
with $r_{\min} = 0$ and $r_{\max} = 0.1$. The noise nudges the sentence
representation into a neighbouring region of the embedding space while
keeping it semantically close to the original. Two augmented copies are
generated per original, tripling the training set. Because the perturbation
is applied to pre-computed embeddings rather than to raw text, no
language-specific vocabulary or grammar rules are needed: the same
procedure works identically for Kinyarwanda, Swahili, and Xitsonga.

\subsubsection{Autoencoder Refinement}
Two autoencoder variants are trained to clean and regularise the
LiDA-perturbed embeddings before they are passed to the classifier.

\textbf{Denoising Autoencoder (DAE).} An encoder
$768 \to 256 \to 64 \to 32$ (LeakyReLU, BatchNorm1d, Dropout 0.2)
compresses the perturbed embedding into a 32-dimensional latent code.
A decoder $32 \to 64 \to 256 \to 768$ reconstructs the original clean
embedding. Training minimises mean-squared reconstruction error at
learning rate $10^{-3}$.

\textbf{Variational Autoencoder (VAE).} A shared encoder layer
($768 \to 512$) branches into two parallel heads producing mean $\mu$
and log-variance $\log\sigma^2$, each of dimension 256. A latent vector
is sampled via the reparameterisation trick and decoded back through
$256 \to 512 \to 768$. The training objective is the ELBO:
\begin{equation}
    \mathcal{L}_{\text{VAE}} =
    \mathbb{E}\!\left[\|\hat{\mathbf{e}} - \mathbf{e}\|_2^2\right]
    + \beta\,D_{\text{KL}}\!\left(
      \mathcal{N}(\mu, \sigma^2)
      \;\big\|\;
      \mathcal{N}(0, I)
    \right)
    \label{eq:vae}
\end{equation}
with $\beta = 1$. The first term penalises reconstruction error (the
autoencoder must reproduce the original embedding faithfully). The second
term penalises deviation from a standard normal, which regularises the
latent space and prevents the model from simply memorising each embedding
without learning general structure. Dropout 0.2 is applied throughout.

\subsubsection{Multi-Head Attention Guided Ensemble (MAGE)}
The MAGE classifier takes a pooled embedding $\mathbf{e}$ and applies
$H = 4$ independent attention heads, each with a learnable context
vector $\mathbf{c}_h \in \mathbb{R}^{768}$:
\begin{equation}
    \mathbf{o}_h =
    \sigma\!\left(\mathbf{e}^{\top} \mathbf{c}_h\right) \cdot \mathbf{e},
    \quad
    \mathbf{z} = \sum_{h=1}^{H} \mathbf{o}_h
    \label{eq:mage}
\end{equation}
The sigmoid acts as a gate: when the embedding aligns well with head $h$'s
context vector, the gate opens and that head contributes strongly to the
final representation $\mathbf{z}$; when alignment is low, the head is
suppressed. Summing four independently-gated copies lets the model
simultaneously attend to four different semantic aspects of the sentence
--- for example, one head may specialise in negation signals while another
attends to intensifiers. The aggregated vector $\mathbf{z}$ is then passed
through a two-layer MLP ($768 \to 384 \to 3$, ReLU, Dropout) to produce
the three sentiment class logits.

\subsection{AfroXLMR-Comet: Knowledge Distillation}\label{sec:distill}

\subsubsection{Teacher and Student Architecture}
The student, AfroXLMR-Comet, is a six-layer XLM-R model with hidden
dimension 256, eight attention heads, and intermediate size 1,024.
Table~\ref{tab:arch} shows the full XLM-R family comparison, including
two intermediate variants (Base and Mini) that serve as upper-bound
baselines in the distillation evaluation.

\begin{table}[htbp]
\caption{XLM-R Model Family: Architecture and Efficiency}
\label{tab:arch}
\begin{center}
\small
\begin{tabular}{lrrrr}
\toprule
\textbf{Property} & \textbf{Large} & \textbf{Base} & \textbf{Mini} & \textbf{Comet} \\
\midrule
Layers            & 24      & 12      & 6      & 6      \\
Hidden size       & 1,024   & 768     & 384    & 256    \\
Attention heads   & 16      & 12      & 12     & 8      \\
Intermediate      & 4,096   & 3,072   & 1,536  & 1,024  \\
\midrule
Parameters        & 559.9M  & 278.0M  & 117.6M & 68.9M  \\
Size (MB)         & 2,135.9 & 1,060.7 & 448.8  & 263.0  \\
Inference (ms)    & 293.9   & 100.5   & 30.2   & 14.0   \\
\bottomrule
\end{tabular}
\end{center}
\end{table}

\subsubsection{Hybrid Distillation Loss}
Training the student with a hybrid loss that combines two complementary
signals produces better results than either signal alone.

The \emph{response-based} component uses KL divergence between softened
teacher and student output distributions:
\begin{equation}
    \mathcal{L}_{\text{distill}} = T^2 \cdot D_{\text{KL}}\!\left(
      \mathrm{softmax}\!\left(\tfrac{\mathbf{z}_s}{T}\right)
      \;\Big\|\;
      \mathrm{softmax}\!\left(\tfrac{\mathbf{z}_t}{T}\right)
    \right)
    \label{eq:kl}
\end{equation}
where $\mathbf{z}_s$ and $\mathbf{z}_t$ are student and teacher logit
vectors and $T = 2$. Dividing by $T$ before the softmax softens the
distributions: instead of the teacher placing near-all probability on the
correct class, it assigns meaningful probability to related classes too,
revealing subtle patterns the student can learn from. The $T^2$ factor
rescales gradient magnitudes to be comparable to a standard cross-entropy.

The \emph{feature-based} component aligns the student's last-layer
attention maps with the teacher's:
\begin{equation}
    \mathcal{L}_{\text{attn}} =
    \left\| \mathbf{W}\,\overline{\mathbf{A}}_s
           - \overline{\mathbf{A}}_t \right\|_2^2
    \label{eq:attn}
\end{equation}
where $\overline{\mathbf{A}} \in \mathbb{R}^{L \times L}$ is the
attention matrix averaged across all heads, then flattened, and
$\mathbf{W}$ is a learnable linear projection mapping the flattened
student attention to the teacher's dimensional space. This forces the
student to develop similar internal attention patterns --- attending to
the same parts of the input --- rather than just producing similar output
labels. The two components are combined with equal weight:
\begin{equation}
    \mathcal{L}_{\text{total}} =
    \alpha\,\mathcal{L}_{\text{distill}}
    + (1 - \alpha)\,\mathcal{L}_{\text{attn}},
    \quad \alpha = 0.5
    \label{eq:total}
\end{equation}

\subsubsection{Training Protocol}
\textbf{Stage 1 --- Distillation.} The teacher is frozen. The student is
trained with AdamW at $\text{lr} = 5 \times 10^{-5}$ for up to 15 epochs,
batch size 8 with 2-step gradient accumulation (effective batch 16), FP16
mixed precision, temperature $T = 2$, early stopping patience 3 and
threshold 0.01. The distillation corpus uses up to 250,000 sequences per
language (Table~\ref{tab:corpus}).

\textbf{Stage 2 --- Fine-tuning.} The distilled student is fine-tuned
independently on each AfriSenti language at
$\text{lr} = 2 \times 10^{-5}$ for 5 epochs. Separating distillation
from fine-tuning prevents task-specific adaptation from interfering with
knowledge transfer.

\begin{table}[htbp]
\caption{Distillation Pretraining Corpus (Stage 1)}
\label{tab:corpus}
\begin{center}
\begin{tabular}{lr}
\toprule
\textbf{Language} & \textbf{Sequences Used} \\
\midrule
Kinyarwanda & 226,466 \\
Swahili     & 250,000 (capped from 537,847) \\
Hausa       & 173,485 \\
Igbo        &  54,410 \\
Yoruba      &  52,067 \\
\bottomrule
\end{tabular}
\end{center}
\end{table}

% ---------------------------------------------------------------
\section{Experiments and Results}\label{sec:results}

\subsection{Cross-Lingual Transfer}\label{sec:res_xlt}

Table~\ref{tab:xlt} reports accuracy and weighted F1 on the Kirundi test
set before and after fine-tuning. All three transformers transfer
strongly: AfriBERTa's zero-shot accuracy (74.21\%) almost matches
mBERT's after fine-tuning (84.62\%), illustrating the value of
Africa-specific pretraining vocabulary. Among traditional models, BiGRU
reaches 83.32\% accuracy after fine-tuning --- competitive with the
transformers --- but its zero-shot accuracy (24.04\%) shows that it has
no transferable knowledge before direct target-language exposure.
CharCNN's post fine-tuning accuracy (48.79\%) confirms that
character-level features alone are insufficient for Bantu morphological
complexity.

\begin{table}[htbp]
\caption{Cross-Lingual Transfer Results on Kirundi Test Set}
\label{tab:xlt}
\begin{center}
\small
\begin{tabular}{lcccc}
\toprule
\multirow{2}{*}{\textbf{Model}} &
\multicolumn{2}{c}{\textbf{Zero-Shot}} &
\multicolumn{2}{c}{\textbf{After Fine-Tune}} \\
\cmidrule(lr){2-3}\cmidrule(lr){4-5}
& Acc & F1 & Acc & F1 \\
\midrule
AfriBERTa  & 0.7421 & 0.7474 & \textbf{0.8830} & \textbf{0.8787} \\
BantuBERTa & 0.7454 & 0.7375 & 0.8657          & 0.8606          \\
mBERT      & 0.5872 & 0.5917 & 0.8462          & 0.8422          \\
\midrule
BiGRU      & 0.2404 & 0.2300 & 0.8332          & 0.8790          \\
TextCNN    & 0.2190 & 0.2320 & 0.5913          & 0.5732          \\
CharCNN    & 0.1916 & 0.1621 & 0.4879          & 0.4764          \\
\bottomrule
\end{tabular}
\end{center}
\end{table}

Table~\ref{tab:forgetting} shows catastrophic forgetting results.
The contrast is stark: mBERT loses only 3.03\% of Kinyarwanda accuracy
after Kirundi fine-tuning; AfriBERTa loses 5.14\%. The traditional
neural models lose 71--75\% --- their Kinyarwanda accuracy collapses
to near-chance level. The most surprising result is BantuBERTa: despite
being an Africa-specific model, it forgets at 74.00\%, on par with the
traditional models. This strongly suggests that \emph{breadth} of
multilingual pretraining, not African-language specificity alone,
determines a model's resistance to catastrophic forgetting.

\begin{table}[htbp]
\caption{Catastrophic Forgetting on Source Language (Kinyarwanda)}
\label{tab:forgetting}
\begin{center}
\begin{tabular}{lccr}
\toprule
\textbf{Model} & \textbf{Acc Before} & \textbf{Acc After} & \textbf{CF\%} \\
\midrule
mBERT          & 0.7884 & 0.7645 & \phantom{0}3.03 \\
AfriBERTa      & 0.8498 & 0.8061 & \phantom{0}5.14 \\
BantuBERTa     & 0.8601 & 0.2172 & 74.00 \\
\midrule
BiGRU          & 0.8851 & 0.2329 & 73.68 \\
TextCNN        & 0.8740 & 0.2207 & 74.86 \\
CharCNN        & 0.6930 & 0.1968 & 71.50 \\
\bottomrule
\multicolumn{4}{l}{\small Acc Before/After: Kinyarwanda accuracy before/after Kirundi fine-tuning.} \\
\end{tabular}
\end{center}
\end{table}

\subsection{MAGE Augmentation Results}\label{sec:res_mage}

Table~\ref{tab:mage} shows results on the combined AfriSenti test set
(Kinyarwanda + Swahili + Xitsonga, 1,482 examples in total). Adding DAE
or VAE alone to the LSTM or logistic regression baseline provides only
marginal gains --- 0.2--1.1 percentage points of accuracy. The MAGE
attention ensemble produces the main improvement: MAGE+DAE improves
accuracy by 3.64 percentage points over the LR baseline; MAGE+VAE
improves by 2.51 points. The attention mechanism contributes +1.21 points
on top of DAE refinement alone. DAE consistently outperforms VAE,
likely because its deterministic reconstruction target is easier to
optimise reliably with only $\sim$7,940 combined training examples.

\begin{table}[htbp]
\caption{MAGE Pipeline Results on Combined AfriSenti Test Set}
\label{tab:mage}
\begin{center}
\small
\begin{tabular}{llcccc}
\toprule
\textbf{Clf.} & \textbf{Aug.} & \textbf{Acc} & \textbf{Prec} & \textbf{Recall} & \textbf{F1} \\
\midrule
LSTM & None  & 0.5680 & 0.5589 & 0.5554 & 0.5566 \\
LSTM & + DAE & 0.5739 & 0.5659 & 0.5697 & 0.5672 \\
LSTM & + VAE & 0.5769 & 0.5681 & 0.5562 & 0.5589 \\
\midrule
LR   & None  & 0.5710 & 0.5636 & 0.5596 & 0.5592 \\
LR   & + DAE & 0.5729 & 0.5672 & 0.5624 & 0.5619 \\
LR   & + VAE & 0.5735 & 0.5659 & 0.5621 & 0.5618 \\
\midrule
MAGE & + DAE & 0.6074\textsuperscript{*} & --- & --- & --- \\
MAGE & + VAE & 0.5961\textsuperscript{*} & --- & --- & --- \\
\bottomrule
\multicolumn{6}{l}{\small \textsuperscript{*}Computed: LR baseline $+$3.64 pp and $+$2.51 pp respectively \cite{b_mage}.} \\
\end{tabular}
\end{center}
\end{table}

\subsection{Knowledge Distillation Results}\label{sec:res_distill}

Table~\ref{tab:distill} compares weighted F1 across five languages.
AfroXLMR-Comet averages 65.28\% F1 against the teacher's 73.22\%,
retaining 89.2\% of teacher performance while being 87.7\% smaller in
parameter count and 21$\times$ faster at inference.
Hausa and Igbo show the strongest retention (71.89\% and 73.82\%)
because their larger fine-tuning corpora (14,172 and 10,192 samples)
give the student more task-specific signal. Swahili shows the weakest
retention (55.30\%) due to its small fine-tuning set (1,810 samples).
Notably, Comet outperforms Mini on Yoruba (66.39\% vs.\ 63.10\%)
despite having 41\% fewer parameters, confirming that the hybrid
distillation loss provides richer initialisation than architectural
compression alone.

\begin{table}[htbp]
\caption{Weighted F1 (\%) Across Languages --- Distillation Results}
\label{tab:distill}
\begin{center}
\small
\begin{tabular}{lcccc}
\toprule
\textbf{Language} & \textbf{Teacher} & \textbf{Base} & \textbf{Mini} & \textbf{Comet} \\
\midrule
Kinyarwanda  & 68.91 & 62.54 & 60.23 & 58.94 \\
Swahili      & 64.10 & 60.81 & 61.88 & 55.30 \\
Hausa        & 79.30 & 78.27 & 74.29 & 71.89 \\
Igbo         & 79.60 & 78.63 & 75.20 & 73.82 \\
Yoruba       & 74.20 & 70.25 & 63.10 & 66.39 \\
\midrule
\textbf{Avg} & 73.22 & 70.10 & 66.94 & 65.28 \\
\bottomrule
\multicolumn{5}{l}{\small Params: Teacher 559.9M, Base 278.0M, Mini 117.6M, Comet 68.9M.} \\
\multicolumn{5}{l}{\small Latency: Teacher 293.9 ms, Base 100.5 ms, Mini 30.2 ms, Comet 14.0 ms.} \\
\end{tabular}
\end{center}
\end{table}

% ---------------------------------------------------------------
\section{Discussion}\label{sec:discussion}

\textbf{Why does AfriBERTa outperform mBERT on cross-lingual transfer?}
AfriBERTa was pretrained exclusively on corpora from 11 African languages,
so its sub-word vocabulary overlaps far more with Kinyarwanda and Kirundi
than mBERT's vocabulary, which must cover 104 languages within the same
token budget. More shared sub-word tokens mean fewer inputs are mapped to
unknown-token fallbacks, and the model's internal representations are
already calibrated to Bantu morphology from pretraining. This explains
both the higher zero-shot transfer (74.21\% vs.\ 58.72\%) and the lower
forgetting rate.

\textbf{Why does BantuBERTa forget catastrophically despite being Africa-specific?}
BantuBERTa was pretrained on a narrower range of Bantu corpora, giving it
a smaller and more specialised vocabulary than AfriBERTa. When fine-tuned
on Kirundi, its parameters shift heavily toward the Kirundi task
distribution, overwriting the Kinyarwanda signal. mBERT's broad 104-language
exposure distributes internal representations across many languages,
making any single fine-tuning step a small perturbation rather than a
complete overwrite. Breadth of multilingual exposure matters more than
African-language specificity for forgetting resilience.

\textbf{Why does MAGE benefit more from DAE than VAE?}
The DAE solves a well-defined reconstruction problem: given a perturbed
embedding, reproduce the clean original. With only $\sim$7,940 training
examples, this converges reliably. The VAE must simultaneously learn a
posterior distribution over the latent space and a reconstruction function
--- a harder optimisation on small data. The VAE's stochastic sampling
also introduces additional variance at test time, which hurts precision.

\textbf{Why does Comet outperform Mini on Yoruba but not Swahili?}
For Yoruba (8,522 fine-tuning samples), the distillation loss gives Comet
a richer initialisation than Mini's random start, and the larger dataset
is enough for Comet to realise that advantage. For Swahili (1,810 samples),
Mini's larger capacity (117.6M vs.\ 68.9M parameters) outweighs the
distillation benefit: there is not enough fine-tuning data for Comet to
translate its better initialisation into higher downstream accuracy.

\textbf{Practical deployment guide.}
The results point to a clear deployment path under different resource
constraints: (1) use AfriBERTa when accuracy matters most and a
126M-parameter model can be served; (2) use AfroXLMR-Comet when latency
($14$~ms) or model size ($263$~MB) is the binding constraint; (3) use
MAGE augmentation when the labelled dataset is very small and no GPU is
available for transformer fine-tuning.

% ---------------------------------------------------------------
\section{Conclusion}\label{sec:conclusion}

We have presented a unified study of three complementary approaches to
low-resource African NLP. Across benchmarks spanning news classification
and tweet sentiment analysis in six African languages, our findings are:

\begin{itemize}
  \item AfriBERTa achieves 88.30\% accuracy on Kirundi (zero-shot:
  74.21\%) and loses only 5.14\% of Kinyarwanda accuracy after target
  fine-tuning, making it the best-performing and most forgetting-resilient
  model for Kinyarwanda--Kirundi transfer.

  \item The CF\% metric reveals an unexpected pattern: BantuBERTa
  (Africa-specific) forgets at 74.00\%, matching the traditional neural
  models (71--75\%), while mBERT (generic multilingual) retains knowledge
  almost perfectly (3.03\%). Breadth of multilingual pretraining --- not
  African-language specificity --- determines forgetting resilience.

  \item MAGE with a denoising autoencoder improves combined AfriSenti
  accuracy by 3.64 percentage points over a logistic regression baseline,
  using no language-specific resources whatsoever. The attention ensemble
  contributes +1.21 points of that improvement beyond the augmented
  embeddings alone.

  \item AfroXLMR-Comet compresses AfroXLMR-Large by 87.7\% in parameter
  count and runs 21$\times$ faster, retaining 89.2\% of teacher F1 across
  five languages. It outperforms the larger AfroXLMR-Mini on Yoruba,
  confirming that hybrid distillation provides richer initialisation than
  architectural compression alone.
\end{itemize}

Future work will extend the CF\% analysis to named entity recognition on
MasakhaNER \cite{b_adelani_ner}, combine MAGE augmentation with
distillation for very-low-resource settings, and evaluate AfroXLMR-Comet
on CPU-only edge hardware representative of Sub-Saharan Africa's
infrastructure.

% ---------------------------------------------------------------
\section*{Acknowledgment}
This work is conducted in collaboration with the Data Science for Social
Impact (DSFSI) Research Laboratory at the University of Pretoria.
We thank the Masakhane NLP community for open-sourcing African language
datasets and pretrained models, and all dataset contributors whose
annotation work made this research possible.

% ---------------------------------------------------------------
\begin{thebibliography}{00}

\bibitem{b_devlin}
J.\ Devlin, M.-W.\ Chang, K.\ Lee, and K.\ Toutanova,
``BERT: Pre-training of deep bidirectional transformers for language understanding,''
in \emph{Proc.\ NAACL-HLT}, 2019, pp.\ 4171--4186.

\bibitem{b_conneau}
A.\ Conneau et al.,
``Unsupervised cross-lingual representation learning at scale,''
in \emph{Proc.\ ACL}, 2020, pp.\ 8440--8451.

\bibitem{b_ogueji}
K.\ Ogueji, Y.\ Zhu, and J.\ Lin,
``Small data? No problem! Exploring the viability of pretrained
multilingual language models for low-resourced languages,''
in \emph{Proc.\ MultilingualNLP Workshop, ACL}, 2021.

\bibitem{b_alabi}
J.\ O.\ Alabi, D.\ Adelani, P.\ Arnaud, and V.\ Marivate,
``Adapting pretrained language models to African languages via
multilingual adaptive fine-tuning,''
in \emph{Proc.\ COLING}, 2022.

\bibitem{b_dossou}
B.\ F.\ P.\ Dossou et al.,
``BantuBERTa: Leveraging multilingual pretrained language models for
Bantu languages,''
in \emph{Proc.\ AfricaNLP Workshop, EACL}, 2022.

\bibitem{b_adelani_afrisenti}
S.\ H.\ Muhammad et al.,
``AfriSenti: A Twitter sentiment analysis benchmark for African languages,''
in \emph{Proc.\ EMNLP}, 2023.

\bibitem{b_kinnews}
R.\ A.\ Niyongabo, H.\ Qu, J.\ Li, and C.\ Ni,
``KINNEWS and KIRNEWS: Benchmarking cross-lingual text classification
for Kinyarwanda and Kirundi,''
in \emph{Proc.\ COLING}, 2020, pp.\ 5507--5521.

\bibitem{b_xlt2024}
``Cross-lingual transfer of multilingual models on low resource African
languages,''
\emph{arXiv:2409.10965}, 2024.

\bibitem{b_mage}
``MAGE: Multi-head attention guided embeddings for low resource sentiment
classification,''
\emph{arXiv:2502.17987}, 2025.

\bibitem{b_comet}
``AfroXLMR-Comet: Multilingual knowledge distillation with attention
matching for low-resource languages,''
\emph{arXiv:2502.18020}, 2025.

\bibitem{b_hinton}
G.\ Hinton, O.\ Vinyals, and J.\ Dean,
``Distilling the knowledge in a neural network,''
\emph{arXiv:1503.02531}, 2015.

\bibitem{b_sanh}
V.\ Sanh, L.\ Debut, J.\ Chaumond, and T.\ Wolf,
``DistilBERT, a distilled version of BERT: smaller, faster, cheaper
and lighter,''
\emph{arXiv:1910.01108}, 2019.

\bibitem{b_jiao}
X.\ Jiao et al.,
``TinyBERT: Distilling BERT for natural language understanding,''
in \emph{Proc.\ EMNLP (Findings)}, 2020, pp.\ 4163--4174.

\bibitem{b_zagoruyko}
S.\ Zagoruyko and N.\ Komodakis,
``Paying more attention to attention: Improving the performance of
convolutional neural networks via attention transfer,''
in \emph{Proc.\ ICLR}, 2017.

\bibitem{b_mccloskey}
M.\ McCloskey and N.\ J.\ Cohen,
``Catastrophic interference in connectionist networks: The sequential
learning problem,''
\emph{Psychol.\ Learning Motivation}, vol.\ 24, pp.\ 109--165, 1989.

\bibitem{b_kirkpatrick}
J.\ Kirkpatrick et al.,
``Overcoming catastrophic forgetting in neural networks,''
\emph{Proc.\ Natl.\ Acad.\ Sci.}, vol.\ 114, no.\ 13,
pp.\ 3521--3526, 2017.

\bibitem{b_wei}
J.\ Wei and K.\ Zou,
``EDA: Easy data augmentation techniques for boosting performance on
text classification tasks,''
in \emph{Proc.\ EMNLP-IJCNLP}, 2019, pp.\ 6382--6388.

\bibitem{b_sennrich}
R.\ Sennrich, B.\ Haddow, and A.\ Birch,
``Improving neural machine translation models with monolingual data,''
in \emph{Proc.\ ACL}, 2016, pp.\ 86--96.

\bibitem{b_kingma}
D.\ P.\ Kingma and M.\ Welling,
``Auto-encoding variational Bayes,''
in \emph{Proc.\ ICLR}, 2014.

\bibitem{b_vincent}
P.\ Vincent, H.\ Larochelle, Y.\ Bengio, and P.-A.\ Manzagol,
``Extracting and composing robust features with denoising autoencoders,''
in \emph{Proc.\ ICML}, 2008, pp.\ 1096--1103.

\bibitem{b_kim}
Y.\ Kim,
``Convolutional neural networks for sentence classification,''
in \emph{Proc.\ EMNLP}, 2014, pp.\ 1746--1751.

\bibitem{b_zhang}
X.\ Zhang, J.\ Zhao, and Y.\ LeCun,
``Character-level convolutional networks for text classification,''
in \emph{Proc.\ NeurIPS}, 2015, pp.\ 649--657.

\bibitem{b_cho}
K.\ Cho et al.,
``Learning phrase representations using RNN encoder-decoder for
statistical machine translation,''
in \emph{Proc.\ EMNLP}, 2014, pp.\ 1724--1734.

\bibitem{b_mikolov}
T.\ Mikolov, I.\ Sutskever, K.\ Chen, G.\ S.\ Corrado, and J.\ Dean,
``Distributed representations of words and phrases and their
compositionality,''
in \emph{Proc.\ NeurIPS}, 2013, pp.\ 3111--3119.

\bibitem{b_vaswani}
A.\ Vaswani et al.,
``Attention is all you need,''
in \emph{Proc.\ NeurIPS}, 2017, pp.\ 5998--6008.

\bibitem{b_wolf}
T.\ Wolf et al.,
``HuggingFace's Transformers: State-of-the-art natural language processing,''
in \emph{Proc.\ EMNLP (System Demonstrations)}, 2020, pp.\ 38--45.

\bibitem{b_adelani_ner}
D.\ Adelani et al.,
``MasakhaNER: Named entity recognition for African languages,''
\emph{Trans.\ Assoc.\ Comput.\ Linguistics}, vol.\ 9,
pp.\ 1116--1131, 2021.

\end{thebibliography}

\end{document}
